\section{Giới hạn và hướng phát triển}
\label{sec:gioi-han}

\subsection{Giới hạn}

\begin{itemize}
  \item \textbf{Phạm vi kênh và hạ tầng}: chỉ hỗ trợ Telegram; một tiến trình,
        SQLite, scheduler in-process --- phù hợp quy mô cá nhân, chưa phân tán.
        Agent loop đồng bộ chiếm một thread cho mỗi tác vụ đang chạy.
  \item \textbf{Bảo mật}: REST API mới bảo vệ bằng một API key đơn (header),
        chưa có auth thật; định danh người dùng dựa vào Telegram.
  \item \textbf{Chống prompt injection ở mức giảm thiểu}: đóng khung dữ liệu,
        vô hiệu hóa delimiter giả, quy tắc trong system prompt, xác nhận người
        dùng --- không có bảo đảm tuyệt đối trước tấn công tinh vi hơn (bài
        toán mở của cả lĩnh vực~\cite{greshake2023injection}).
  \item \textbf{Trần chi phí theo chặng chạy}: trần 0{,}25~USD áp cho từng
        chặng run/resume, không cộng dồn qua các lần tạm dừng chờ xác nhận ---
        một tác vụ nhiều lần confirm có thể tốn hơn trần danh nghĩa (tổng chi
        phí thật vẫn tra được qua \texttt{task\_usage}).
  \item \textbf{Rate limit và cache in-memory}: trạng thái rate limiter và
        search cache nằm trong RAM của tiến trình --- restart là reset; cache
        tìm kiếm khớp theo truy vấn chuẩn hóa (exact-match), chưa phải
        semantic cache theo nghĩa embedding.
  \item \textbf{Memory dài hạn}: mới dừng ở hồ sơ người dùng dạng JSON, agent
        chưa khai thác; chưa có vector memory/RAG.
  \item \textbf{Giới hạn kỹ thuật}: timeout không kill được thread công cụ
        đang chạy (giới hạn Python) --- thread cũ được bỏ lại có kiểm soát.
  \item \textbf{Đánh giá}: LLM-as-judge mới chấm pass/fail theo tiêu chí từng
        case, chưa có thang điểm nhiều mức và chưa kiểm chứng chéo tay trên
        mẫu lớn; trong lần chạy mẫu, judge còn dùng chính model agent
        (\texttt{gemini-3.1-flash-lite}) nên có rủi ro self-preference.
        Bộ 66 case chưa phủ hết không gian hành vi.
  \item \textbf{Phần model thật chỉ là mẫu minh họa}: quota free tier Gemini
        (500 request/ngày) cạn giữa chừng nên thay vì đủ ma trận
        $2 \times 2$, chỉ đo được một cấu hình ReAct $\times$ Gemini trên
        22/66 case ($n=3$, 66 run) --- chưa đủ để so sánh chéo model hay so
        sánh hai chiến lược trên model thật (RQ1, RQ2 mới trả lời một
        phần, Mục~\ref{sec:ket-qua}); việc đo model khác sẽ thực hiện trong
        tương lai. Số liệu exp-final cũng chỉ ra hai điểm cần thận trọng:
        (1) kết quả model thật dao động giữa các run (2/22 case đổi
        pass/fail giữa 3 lượt) nên khác biệt nhỏ giữa cấu hình chưa có ý
        nghĩa; (2) nhóm yếu nhất trên model thật là multi-step
        (52{,}4\%) và phục hồi lỗi (0/3 --- model báo lỗi trung thực nhưng
        không tự đổi nguồn); độ trễ đo được còn bị méo bởi backoff 429 của
        free tier (trung bình 25{,}7\,s so với trung vị 4{,}9\,s).
\end{itemize}

\subsection{Hướng phát triển}

\begin{itemize}
  \item Mở rộng kênh (web chat/Slack) trên cùng lõi agent; PostgreSQL và
        scheduler phân tán khi cần đa máy; tách worker + queue để bỏ giới hạn
        một thread/tác vụ.
  \item Memory dài hạn có cấu trúc + RAG cho tri thức cá nhân.
  \item Phòng thủ injection nhiều lớp hơn (bộ lọc đầu ra công cụ, sandbox
        quyền theo từng tác vụ); trần chi phí cộng dồn theo cả vòng đời tác
        vụ.
  \item Chuẩn hóa tool interface theo hướng Model Context Protocol (MCP) để
        dùng lại hệ sinh thái công cụ bên ngoài.
  \item Mở rộng eval: chạy đủ ma trận 2 chiến lược $\times$ nhiều model thật
        (khi có quota hoặc model khác) để trả lời trọn RQ1--RQ2; thêm case
        đối kháng, tăng số lượt lặp, thang điểm judge nhiều mức có kiểm
        chứng người (judge tách model với agent), đo chi phí--chất lượng
        theo Pareto để chọn model.
\end{itemize}

\section*{Kết luận}
\addcontentsline{toc}{section}{Kết luận}

Đồ án xây dựng Laplace's Demon --- một AI agent cá nhân hoàn chỉnh từ kênh
giao tiếp (Telegram) đến lõi điều phối (state machine tự xây với hai chiến
lược ReAct và Plan-and-Execute cắm chung một interface), với ba trụ cột thiết
kế được hiện thực hóa nhất quán: an toàn theo lớp (giới hạn bước/timeout,
xác nhận người dùng bền vững qua restart, trần chi phí, rate limit, chống
prompt injection nhiều lớp), khả năng quan sát (mọi lệnh gọi LLM và mọi lần
chạy công cụ đều có trace; trace viewer kèm chi phí và chế độ phát lại
offline), và đánh giá định lượng (eval harness 66 case, 8 metric, chạy đúng
agent loop production, tùy chọn LLM-as-judge). Thực nghiệm
\texttt{exp-final} cho thấy khung agent đạt 100\% trên toàn bộ 198 run
mock trùng chiến lược gốc ở cả hai chiến lược (vai trò regression
benchmark), và trên mẫu model thật Gemini $\times$ ReAct (22 case,
$n=3$) đạt 72{,}7\% success với chi phí $\approx$ 0{,}0005 USD/tác vụ ---
mạnh ở phân tuyến (route accuracy 94{,}4\%), trả lời trực tiếp và chống
injection (không run nào bị vượt), yếu ở multi-step và tự phục hồi lỗi;
ReAct được giữ làm chiến lược mặc định vì là cấu hình duy nhất đã có số
liệu model thật và xử lý mềm hơn khi người dùng từ chối confirm.

Bài học lớn nhất của đồ án nằm ở phương pháp: khi trace là công dân hạng
nhất và bộ eval lặp lại được, mọi thay đổi --- từ sửa prompt đến đổi chiến
lược hay model --- đều trở thành một thí nghiệm có số liệu thay vì một cảm
nhận; chính vòng phản hồi này (eval phát hiện model bịa tên công cụ $\to$
sửa prompt $\to$ eval xác nhận cải thiện) đã định hình phần lớn các kỹ thuật
ở Mục~\ref{sec:ky-thuat}.
